\documentclass[11pt]{article}

\usepackage[margin=1in]{geometry}
\usepackage{parskip}
\usepackage{amsmath, amssymb, amsthm}
\usepackage{booktabs}
\usepackage{hyperref}

\title{A small library for finite discrete-time Markov chains}
\author{Group N}
\date{}

\begin{document}
\maketitle

\begin{abstract}
We present \texttt{markov}, a small Python library for finite
discrete-time Markov chains. The library represents a chain by its
transition matrix as a \texttt{sympy.Matrix}~\cite{meurer2017} and
exposes three short functions: a stochasticity check, the distribution
after \(n\) steps, and the stationary distribution. We illustrate the
library on a two-state chain whose stationary distribution we compute
exactly as \((4/7, 3/7)\), and we verify numerically that the
distribution after \(50\) steps from a deterministic start agrees with
the stationary distribution to within \(10^{-10}\). The library follows
the Diataxis framework for documentation~\cite{procida2017} and the
matrices chapter of the open-source textbook \emph{Python for
Mathematics}~\cite{knight2024}.
\end{abstract}

\section{Introduction}

A discrete-time Markov chain on a finite state space is a stochastic
process where the next state depends only on the current state and not
on the history. The chain is described by its transition matrix \(P\),
a row-stochastic matrix where \(P_{ij}\) is the probability of moving
from state \(i\) to state \(j\) in one step~\cite{norris1998}. The
distribution of the chain after \(n\) steps, starting from the row
vector \(\pi_0\) of initial probabilities, is

\[
    \pi_n = \pi_0 P^n.
\]

A stationary distribution is a row vector \(\pi\) such that
\(\pi P = \pi\). For an irreducible, aperiodic chain on a finite state
space, the stationary distribution exists and is unique, and the
distribution after \(n\) steps converges to it from any starting
point~\cite{norris1998}.

We wanted a small library that we could use to inspect the
distributions of a chain step by step and to compute the stationary
distribution exactly. The \texttt{numpy} package can multiply matrices
quickly but loses the exact rational form of the answer. We chose to
build on \texttt{sympy.Matrix} instead, following the conventions of
the matrices chapter of \emph{Python for
Mathematics}~\cite{knight2024}.

We present a three-function library,
\texttt{markov}; we recover the stationary distribution of a
two-state chain in closed form; and we show numerically that \(P^n\)
applied to a deterministic initial distribution agrees with the
stationary distribution to within \(10^{-10}\) after \(50\) steps.

\section{The library}

The library lives in the single module \texttt{markov.py} and exposes
three functions.

\texttt{markov.is\_stochastic(transition\_matrix)}
returns \texttt{True} if every row of the matrix sums to one. We use
\texttt{sympy.simplify} on the difference so that the check works even
when the entries are symbolic.

\texttt{markov.n\_step\_distribution(initial\_distribution,
transition\_matrix, number\_of\_steps)} multiplies the row vector
\(\pi_0\) by \(P^n\). The matrix power is computed by \texttt{sympy}
using fast exponentiation, so even moderately large values of \(n\)
are cheap.

\texttt{markov.stationary\_distribution(transition\_matrix)} returns
the row vector \(\pi\) with \(\pi P = \pi\) and \(\sum_i \pi_i = 1\).
We compute it as the null space of \(P^T - I\) and normalise. This
direct approach is short and works on any matrix; for very large
chains a sparse implementation would be more efficient, but the
library is designed for small chains where exactness matters more
than speed.

\section{Worked example}

We illustrate the library on the two-state chain with transition matrix

\[
    P = \begin{pmatrix} 7/10 & 3/10 \\ 4/10 & 6/10 \end{pmatrix}.
\]

The stationary distribution can be computed by hand: solving
\(\pi P = \pi\) with \(\pi_0 + \pi_1 = 1\) gives \(\pi = (4/7, 3/7)\).
The library returns this answer exactly. Starting from
\(\pi_0 = (1, 0)\), the distributions after \(1\), \(2\) and \(50\)
steps are shown in Table~\ref{tab:steps}.

\begin{table}[h]
\centering
\begin{tabular}{l l}
\toprule
\(n\) & \(\pi_n\) \\
\midrule
\(0\)  & \((1, 0)\) \\
\(1\)  & \((7/10, 3/10)\) \\
\(2\)  & \((61/100, 39/100)\) \\
\(50\) & numerically \((0.571429, 0.428571)\) \\
\bottomrule
\end{tabular}
\caption{\textbf{Convergence to the stationary distribution.}
Starting from \(\pi_0 = (1, 0)\), the distribution after \(n\) steps
moves quickly towards the stationary distribution \((4/7, 3/7) =
(0.571\dots, 0.428\dots)\). After \(50\) steps the error in each
component is below \(10^{-10}\).}
\label{tab:steps}
\end{table}

\section{Discussion}

For small chains the stationary
distribution often has a clean rational form (such as \((4/7, 3/7)\))
that a numeric library would round to a decimal. We wanted to recover
the closed form, and \texttt{sympy.Matrix} keeps every step exact.

The library handles only finite chains, and only
the discrete-time case. A natural extension is to add continuous-time
chains, where the rate matrix \(Q\) plays the role of \(P - I\) and
the distribution after time \(t\) is \(\pi_0 \exp(Q t)\). Symbolic
matrix exponentials are available in \texttt{sympy} so this would
follow the same shape as the existing code.

\section{Conclusion}

We have presented \texttt{markov}, a three-function library for
finite discrete-time Markov chains, written on top of
\texttt{sympy.Matrix}. The library recovers the stationary
distribution of a two-state chain in closed form and the iterates
converge to it as expected.

\bibliographystyle{plain}
\bibliography{references}

\end{document}
